% ══════════════════════════════════════════════════════════════════════════════
\section{Results}
\label{sec:results}
% ══════════════════════════════════════════════════════════════════════════════

\subsection{Main Results}

Table~\ref{tab:main_results} reports estimates from the preferred 3-month forward conflict window across four specifications: pooled OLS without fixed effects (Column~1), TWFE OLS (Column~2), TWFE Poisson QMLE (Column~3), and the preferred TWFE OLS log-log specification (Column~4).

\begin{table}[H]
\centering
\caption{Effect of Predicted Maize Yield on Conflict: Main Results (3-Month Window)}
\label{tab:main_results}
\begin{threeparttable}

\begingroup
\centering
\begin{tabular}{lcccc}
   \tabularnewline \midrule \midrule
   Dependent Variables: & \multicolumn{3}{c}{Conflict events (3-mo.)} & Log conflict (3-mo.)\\
                           & (1)\\Pooled OLS       & (2)\\TWFE OLS   & (3)\\TWFE Poisson   & (4)\\TWFE OLS Log-Log \\     
   Model:                  & (1)                   & (2)             & (3)                 & (4)\\  
                           &  OLS                  & OLS             & Poisson             & OLS\\  
   \midrule
   \emph{Variables}\\
   Constant                & 1.450$^{***}$         &                 &                     &   \\   
                           & (0.2119)              &                 &                     &   \\   
   Log pred. yield (kg/ha) & -0.6383$^{***}$       & -3.750$^{***}$  & -2.629$^{***}$      & -0.6656$^{***}$\\   
                           & (0.2105)              & (1.155)         & (0.8555)            & (0.1533)\\   
   Log population          &                       & 7.193$^{***}$   & -2.651$^{***}$      & 1.460$^{***}$\\   
                           &                       & (1.890)         & (0.9572)            & (0.2833)\\   
   \midrule
   \emph{Fixed-effects}\\
   Admin-2 FE              &                       & Yes             & Yes                 & Yes\\  
   Year FE                 &                       & Yes             & Yes                 & Yes\\  
   \midrule
   \emph{Fit statistics}\\
   Observations            & 17,889                & 17,883          & 11,100              & 17,883\\  
   Squared Correlation     & 0.00051               & 0.30587         & 0.50175             & 0.47954\\  
   Pseudo R$^2$            & $9.35\times 10^{-5}$  & 0.06644         & 0.58150             & 0.36091\\  
   BIC                     & 98,311.9              & 104,777.8       & 35,374.7            & 33,712.8\\  
   \midrule \midrule
   \multicolumn{5}{l}{\emph{Signif. Codes: ***: 0.01, **: 0.05, *: 0.1}}\\
\end{tabular}
 
\par \raggedright 
\textit{Notes:} Clustered standard errors at the Admin-2 level in parentheses. Column~(1) is pooled OLS with no fixed effects or controls. Columns~(2)--(4) include Admin-2 and year two-way fixed effects. The key regressor is the log of v2 XGBoost-predicted maize yield (kg/ha), trained on GROW-Africa administrative yield records with satellite, soil, and climate inputs (spatial $R^2 = 0.35$). The log population control is log of WorldPop gridded population summed to Admin-2; values for 2021--2024 are carried forward from 2020. Column~(4) is the preferred specification. Uganda is excluded from all analyses (not in model training data; LOCO $R^2 < 0$). $^{*}$ $p<0.10$, $^{**}$ $p<0.05$, $^{***}$ $p<0.01$.
\par\endgroup



\end{threeparttable}
\end{table}

The pooled OLS estimate (Column~1) is negative ($\hat{\beta} = -0.650$, $p = 0.002$), but this cross-sectional association conflates within-district variation with unobserved heterogeneity across countries and regions. The TWFE OLS specification (Column~2), which absorbs district and year fixed effects, yields a negative and statistically significant coefficient of $\hat{\beta} = -3.750$ ($\text{SE} = 1.155$, $p = 0.001$). The large standard error reflects the high variance of raw conflict counts driven by extreme events in a small number of high-conflict district-years.

The TWFE Poisson specification (Column~3) also produces a negative and significant estimate ($\hat{\beta} = -2.629$, $p = 0.002$). The Poisson model is estimated on the 11,100 district-year observations that survive the Poisson absorbing condition, which drops all districts recording zero conflict events in every year of the sample. The coefficient is interpreted as a semi-elasticity: a 1\% increase in predicted yield is associated with a 2.6\% reduction in the expected event count among persistently conflict-affected districts.

The preferred specification is Column~(4), TWFE OLS log-log. The estimated elasticity is $\hat{\beta} = -0.666$ ($\text{SE} = 0.153$, $p < 0.001$): a 10\% increase in ML-predicted maize yield is associated with a 6.7\% reduction in log conflict events in the subsequent 3-month window, conditional on district and year fixed effects and log population. The $\log(1+x)$ transformation of the outcome compresses the right tail of the highly skewed conflict distribution and delivers more precise estimates than the linear specifications while retaining all zero-conflict observations. The within-$R^2$ of 0.014 is modest in absolute terms but consistent with the literature on conflict data, where fixed effects absorb the bulk of explained variation and within-unit identification relies on year-to-year changes in satellite-derived yield predictions \citep{HarariLaFerrara2018}.

The negative and significant yield effect is consistent with the opportunity-cost channel: higher agricultural productivity raises the return to farming relative to participation in or support for armed activity, reducing the incentive to engage in conflict \citep{DalBoDalBo2011}. This interpretation is consistent with the broader empirical literature on income shocks and civil conflict in Africa \citep{MiguelSatyanathSergenti2004, BlattmanMiguel2010}.

\subsection{Robustness Check 1: Alternative Lag Windows}

Table~\ref{tab:robustness} tests whether the negative yield--conflict relationship persists across alternative forward conflict windows. Columns~(1)--(3) replicate the preferred log-log OLS specification for 3-, 6-, and 12-month windows respectively.

\begin{table}[H]
\centering
\caption{Robustness Checks: Alternative Windows and Outcome Definitions}
\label{tab:robustness}
\begin{threeparttable}

\begingroup
\centering
\begin{tabular}{lccccc}
   \tabularnewline \midrule \midrule
   Dependent Variables:    & Log conflict (3-mo.)      & Log conflict (6-mo.)      & Log conflict (12-mo.)      & Conflict events (3-mo.)   & Fatalities (3-mo.)\\  
                           & (1)\\3-Month\\Log-Log     & (2)\\6-Month\\Log-Log     & (3)\\12-Month\\Log-Log     & (4)\\3-Month\\Poisson     & (5)\\Fatalities\\Poisson \\       
   Model:                  & (1)                       & (2)                       & (3)                        & (4)                       & (5)\\  
                           &  OLS                      & OLS                       & OLS                        & Poisson                   & Poisson\\  
   \midrule
   \emph{Variables}\\
   Log pred. yield (kg/ha) & -0.6656$^{***}$           & -0.9488$^{***}$           & -1.290$^{***}$             & -2.629$^{***}$            & 0.2725\\   
                           & (0.1533)                  & (0.1927)                  & (0.2356)                   & (0.8555)                  & (1.325)\\   
   Log population          & 1.460$^{***}$             & 1.739$^{***}$             & 2.002$^{***}$              & -2.651$^{***}$            & -2.996$^{*}$\\   
                           & (0.2833)                  & (0.3455)                  & (0.3908)                   & (0.9572)                  & (1.547)\\   
   \midrule
   \emph{Fixed-effects}\\
   Admin-2 FE              & Yes                       & Yes                       & Yes                        & Yes                       & Yes\\  
   Year FE                 & Yes                       & Yes                       & Yes                        & Yes                       & Yes\\  
   \midrule
   \emph{Fit statistics}\\
   Observations            & 17,883                    & 17,883                    & 17,883                     & 11,100                    & 9,100\\  
   Squared Correlation     & 0.47954                   & 0.54730                   & 0.61484                    & 0.50175                   & 0.27690\\  
   Pseudo R$^2$            & 0.36091                   & 0.34003                   & 0.34082                    & 0.58150                   & 0.53311\\  
   BIC                     & 33,712.8                  & 40,540.1                  & 46,032.4                   & 35,374.7                  & 107,589.3\\  
   \midrule \midrule
   \multicolumn{6}{l}{\emph{Clustered (Admin-2 FE) standard-errors in parentheses}}\\
   \multicolumn{6}{l}{\emph{Signif. Codes: ***: 0.01, **: 0.05, *: 0.1}}\\
\end{tabular}
 
\par \raggedright 
\textit{Notes:} Clustered standard errors at the Admin-2 level in parentheses. All models include Admin-2 and year two-way fixed effects. Columns~(1)--(3) use the log$(1+\text{conflict events})$ outcome under the TWFE OLS log-log specification, varying the length of the forward conflict window (3, 6, and 12 months post-harvest). Column~(4) uses conflict event counts under TWFE Poisson QMLE. Column~(5) replaces the conflict count with fatality count under Poisson QMLE; districts with zero fatalities in all years are dropped by the Poisson absorbing condition. $^{*}$ $p<0.10$, $^{**}$ $p<0.05$, $^{***}$ $p<0.01$.
\par\endgroup



\end{threeparttable}
\end{table}

The negative yield effect strengthens monotonically as the window lengthens: the estimated elasticity rises from $-0.666$ at 3 months to $-0.949$ ($p < 0.001$) at 6 months and $-1.290$ ($p < 0.001$) at 12 months. This pattern is consistent with a mechanism that unfolds gradually over the post-harvest economic cycle rather than operating exclusively at the immediate harvest margin. Higher yields may reduce conflict through several reinforcing channels over the medium term: replenishment of household food stocks and productive assets, reduced vulnerability to armed recruitment during the lean season, and more sustained increases in the opportunity cost of participation in violence.

Column~(4) replicates the Poisson count specification at 3 months; the coefficient ($\hat{\beta} = -2.629$, $p = 0.002$) confirms the negative direction under a count-outcome model \citep{SantosSilvaTenreyro2006}. Column~(5) uses fatality counts as the outcome; the coefficient is positive and insignificant ($\hat{\beta} = 0.272$, $p = 0.837$). This divergence between the event-count and fatality results suggests that higher yields reduce the frequency of conflict incidents rather than their average lethality---a pattern more consistent with the deterrence of small-scale predatory events than with the disruption of large-scale organised campaigns.

\subsection{Robustness Check 2: Spatial Lag Control}

A potential concern is that conflict exhibits strong spatial clustering \citep{HarariLaFerrara2018}, and that uncontrolled spillovers from neighbouring districts could confound the within-district yield--conflict estimate. Table~\ref{tab:spatial_lag} augments the baseline specification with a spatial lag of the conflict outcome, defined as the unweighted mean conflict events among all queen-contiguous neighbouring districts in the same year, computed from the GADM Level-2 polygon adjacency matrix.

\begin{table}[H]
\centering
\caption{Robustness Check: Spatial Lag of Conflict}
\label{tab:spatial_lag}
\begin{threeparttable}

\begingroup
\centering
\begin{tabular}{lcccc}
   \tabularnewline \midrule \midrule
   Dependent Variables: & \multicolumn{2}{c}{Log conflict (3-mo.)} & \multicolumn{2}{c}{Conflict events (3-mo.)}\\
                           & (1)\\OLS\\Baseline     & (2)\\OLS\\Spatial Lag     & (3)\\Poisson\\Baseline     & (4)\\Poisson\\Spatial Lag \\       
   Model:                  & (1)                    & (2)                       & (3)                        & (4)\\  
                           &  OLS                   & OLS                       & Poisson                    & Poisson\\  
   \midrule
   \emph{Variables}\\
   Log pred. yield (kg/ha) & -0.6656$^{***}$        & -0.3765$^{***}$           & -2.629$^{***}$             & -0.9577\\   
                           & (0.1533)               & (0.1097)                  & (0.8555)                   & (0.7474)\\   
   Log population          & 1.460$^{***}$          & 0.9089$^{***}$            & -2.651$^{***}$             & -2.593$^{***}$\\   
                           & (0.2833)               & (0.2469)                  & (0.9572)                   & (0.9765)\\   
   Spatial lag (conflict)  &                        & 0.0932$^{***}$            &                            & 0.0654$^{***}$\\   
                           &                        & (0.0059)                  &                            & (0.0068)\\   
   \midrule
   \emph{Fixed-effects}\\
   Admin-2 FE              & Yes                    & Yes                       & Yes                        & Yes\\  
   Year FE                 & Yes                    & Yes                       & Yes                        & Yes\\  
   \midrule
   \emph{Fit statistics}\\
   Observations            & 17,883                 & 17,883                    & 11,100                     & 11,100\\  
   Squared Correlation     & 0.47954                & 0.59560                   & 0.50175                    & 0.60955\\  
   Pseudo R$^2$            & 0.36091                & 0.50035                   & 0.58150                    & 0.61478\\  
   BIC                     & 33,712.8               & 29,210.5                  & 35,374.7                   & 33,163.7\\  
   \midrule \midrule
   \multicolumn{5}{l}{\emph{Clustered (Admin-2 FE) standard-errors in parentheses}}\\
   \multicolumn{5}{l}{\emph{Signif. Codes: ***: 0.01, **: 0.05, *: 0.1}}\\
\end{tabular}
 
\par \raggedright 
\textit{Notes:} Clustered standard errors at the Admin-2 level in parentheses. All models include Admin-2 and year two-way fixed effects. Columns~(1) and~(3) reproduce the baseline log-log OLS and Poisson specifications without the spatial lag. Columns~(2) and~(4) augment those specifications with a spatial lag of the conflict outcome, defined as the population-weighted mean of conflict events in queen-contiguous neighbouring districts. Contiguity is computed from GADM Level-2 polygons using the \texttt{spdep} package. Districts with no contiguous neighbours receive a spatial lag of zero. $^{*}$ $p<0.10$, $^{**}$ $p<0.05$, $^{***}$ $p<0.01$.
\par\endgroup



\end{threeparttable}
\end{table}

The spatial lag enters with a large and highly significant positive coefficient in both OLS (Column~2: $\hat{\beta}_\text{lag} = 0.093$, $p < 0.001$) and Poisson (Column~4: $\hat{\beta}_\text{lag} = 0.065$, $p < 0.001$) specifications, confirming substantial spatial clustering of conflict: a district surrounded by violent neighbours is significantly more likely to experience conflict itself. Conditioning on the spatial lag substantially improves model fit (within-$R^2$ rises from 0.014 to 0.234 in OLS).

Critically, the yield coefficient remains negative and statistically significant after controlling for spatial spillovers: the OLS log-log estimate falls from $-0.666$ to $-0.377$ ($\text{SE} = 0.110$, $p < 0.001$). The Poisson coefficient also remains negative ($\hat{\beta} = -0.958$) but loses significance ($p = 0.200$). The persistence of a negative and significant yield effect in the OLS specification after conditioning on neighbour conflict indicates that the main result is not an artefact of spatial autocorrelation. The attenuation of the coefficient from $-0.666$ to $-0.377$ suggests that some of the raw yield effect is mediated through spatial contagion---districts with high yields may be shielded from conflict that originates in and spreads from lower-yield neighbours---but a direct within-district yield effect on own-district conflict remains.

\subsection{Heterogeneous Effects 1: Conflict Type}

Table~\ref{tab:conflict_types} disaggregates the analysis by ACLED event type. Each column reports the preferred TWFE OLS log-log specification with a type-specific log$(1+\text{events})$ outcome for the 3-month window.

\begin{table}[H]
\centering
\caption{Heterogeneous Effects by Conflict Type (TWFE OLS Log-Log, 3-Month Window)}
\label{tab:conflict_types}
\begin{threeparttable}

\begingroup
\centering
\begin{tabular}{lcccc}
   \tabularnewline \midrule \midrule
   Dependent Variables:      & Log battles (3-mo.) & Log riots (3-mo.) & Log viol.\ vs.~civilians (3-mo.)  & Log expl./remote (3-mo.)\\  
                             & (1)\\Battles        & (2)\\Riots        & (3)\\Viol.\ vs.\ Civ.             & (4)\\Expl./Remote \\     
   Model:                    & (1)                 & (2)               & (3)                               & (4)\\  
   \midrule
   \emph{Variables}\\
   Log pred.\ yield (kg/ha)  & -0.3086$^{***}$     & -0.0049           & -0.5236$^{***}$                   & -0.1549$^{**}$\\   
                             & (0.0989)            & (0.0348)          & (0.1060)                          & (0.0665)\\   
   Log population            & 0.5983$^{***}$      & 0.1242            & 1.026$^{***}$                     & 0.3645$^{***}$\\   
                             & (0.1636)            & (0.1105)          & (0.1862)                          & (0.1339)\\   
   \midrule
   \emph{Fixed-effects}\\
   Admin-2 FE                & Yes                 & Yes               & Yes                               & Yes\\  
   Year FE                   & Yes                 & Yes               & Yes                               & Yes\\  
   \midrule
   \emph{Fit statistics}\\
   Observations              & 17,883              & 17,883            & 17,883                            & 17,883\\  
   R$^2$                     & 0.37349             & 0.25528           & 0.36094                           & 0.35622\\  
   Within R$^2$              & 0.00550             & 0.00065           & 0.01350                           & 0.00483\\  
   \midrule \midrule
   \multicolumn{5}{l}{\emph{Clustered (Admin-2 FE) standard-errors in parentheses}}\\
   \multicolumn{5}{l}{\emph{Signif. Codes: ***: 0.01, **: 0.05, *: 0.1}}\\
\end{tabular}
 
\par \raggedright 
\textit{Notes:} Clustered standard errors at the Admin-2 level in parentheses. All models are TWFE OLS log-log, estimated on the full seven-country panel. Each column uses a conflict-type-specific outcome: log$(1+\text{events})$ in the 3-month post-harvest window for battles, riots, violence against civilians, and explosions/remote violence respectively. Event-type disaggregation is based on the ACLED \texttt{event\_type} classification. All models include Admin-2 and year two-way fixed effects. $^{*}$ $p<0.10$, $^{**}$ $p<0.05$, $^{***}$ $p<0.01$.
\par\endgroup



\end{threeparttable}
\end{table}

The yield effect differs substantially across event types. It is negative and statistically significant for battles ($\hat{\beta} = -0.309$, $\text{SE} = 0.099$, $p = 0.002$), violence against civilians ($\hat{\beta} = -0.524$, $\text{SE} = 0.106$, $p < 0.001$), and explosions/remote violence ($\hat{\beta} = -0.155$, $\text{SE} = 0.066$, $p = 0.020$). The coefficient for riots is statistically indistinguishable from zero ($\hat{\beta} = -0.005$, $\text{SE} = 0.035$, $p = 0.887$).

The largest effect is for violence against civilians. This category captures deliberate targeting of non-combatants by armed actors---a form of violence whose incidence reflects a calculated decision to expend resources on predatory activity. A reduction in this type of violence when yields are higher is directly consistent with the opportunity-cost channel operating at the armed-actor margin: when agricultural income is higher, the relative return to predatory violence falls, and armed groups that engage in civilian targeting reduce the frequency of such attacks. The significant negative effect on explosions and remote violence---which requires logistical and financial resources to execute---reinforces the interpretation that better harvests reduce the profitability or feasibility of resource-intensive predatory operations.

The null result for riots is theoretically informative. Riot activity reflects collective civilian mobilisation, driven predominantly by grievance, political contention, and urban economic conditions. Rural yield shocks operate primarily through the economic opportunity cost and predation channels, both of which concern armed-group or combatant behaviour rather than civilian collective action. The absence of a yield effect on riots is therefore consistent with the proposed mechanisms and provides implicit evidence against the hypothesis that the yield--conflict relationship is driven by a broad economic discontent channel that would affect all types of conflict indiscriminately.

\subsection{Heterogeneous Effects 2: Agro-Ecological Zone}

Table~\ref{tab:aez} examines heterogeneity in the yield--conflict elasticity between the Sahel and Non-Sahel zones. Columns~(1) and~(2) report split-sample TWFE OLS log-log estimates; Column~(3) presents the full-sample interaction model in which log predicted yield is interacted with a Sahel-zone indicator; Columns~(4) and~(5) report Poisson estimates by zone.

\begin{table}[H]
\centering
\caption{Heterogeneous Effects by Agro-Ecological Zone: Sahel vs.\ Non-Sahel}
\label{tab:aez}
\begin{threeparttable}

\begingroup
\centering
\begin{tabular}{lccccc}
   \tabularnewline \midrule \midrule
   Dependent Variables: & \multicolumn{3}{c}{Log conflict (3-mo.)} & \multicolumn{2}{c}{Conflict events (3-mo.)}\\
                             & (1)\\Non-Sahel\\OLS     & (2)\\Sahel\\OLS     & (3)\\Full Sample\\OLS (Interaction)     & (4)\\Non-Sahel\\Poisson     & (5)\\Sahel\\Poisson \\       
   Model:                    & (1)                     & (2)                 & (3)                                     & (4)                         & (5)\\  
                             &  OLS                    & OLS                 & OLS                                     & Poisson                     & Poisson\\  
   \midrule
   \emph{Variables}\\
   Log pred.\ yield (kg/ha)  & -0.4666$^{***}$         & -0.8769             & -1.130$^{***}$                          & 1.635                       & -3.950$^{***}$\\   
                             & (0.1227)                & (0.6984)            & (0.1704)                                & (1.153)                     & (1.094)\\   
   Log population            & 0.9731$^{***}$          & -0.3672             & 0.8516$^{***}$                          & -1.572                      & -6.228$^{***}$\\   
                             & (0.2255)                & (0.7136)            & (0.2706)                                & (1.066)                     & (1.985)\\   
   Log yield $\times$ Sahel  &                         &                     & 3.513$^{***}$                           &                             &   \\   
                             &                         &                     & (0.4166)                                &                             &   \\   
   \midrule
   \emph{Fixed-effects}\\
   Admin-2 FE                & Yes                     & Yes                 & Yes                                     & Yes                         & Yes\\  
   Year FE                   & Yes                     & Yes                 & Yes                                     & Yes                         & Yes\\  
   \midrule
   \emph{Fit statistics}\\
   Observations              & 16,017                  & 1,866               & 17,883                                  & 9,294                       & 1,806\\  
   Squared Correlation       & 0.43582                 & 0.59466             & 0.48971                                 & 0.36979                     & 0.69338\\  
   Pseudo R$^2$              & 0.39424                 & 0.31091             & 0.37181                                 & 0.50351                     & 0.69179\\  
   BIC                       & 25,753.2                & 4,804.2             & 33,369.8                                & 26,418.0                    & 7,476.7\\  
   \midrule \midrule
   \multicolumn{6}{l}{\emph{Clustered (Admin-2 FE) standard-errors in parentheses}}\\
   \multicolumn{6}{l}{\emph{Signif. Codes: ***: 0.01, **: 0.05, *: 0.1}}\\
\end{tabular}
 
\par \raggedright 
\textit{Notes:} Clustered standard errors at the Admin-2 level in parentheses. All models include Admin-2 and year two-way fixed effects. The Sahel zone comprises Burkina Faso, Mali, and Niger; the Non-Sahel zone comprises Ethiopia, Malawi, Nigeria, and Tanzania. Columns~(1) and~(2) are split-sample TWFE OLS log-log specifications. Column~(3) is the full-sample interaction model; the coefficient on \textit{Log yield $\times$ Sahel} gives the differential yield effect for the Sahel zone relative to the Non-Sahel baseline: the net Sahel effect is the sum of \textit{Log pred.\ yield} and \textit{Log yield $\times$ Sahel}. Columns~(4) and~(5) repeat the split-sample estimation under Poisson QMLE. $^{*}$ $p<0.10$, $^{**}$ $p<0.05$, $^{***}$ $p<0.01$.
\par\endgroup



\end{threeparttable}
\end{table}

The Non-Sahel split-sample OLS estimate (Column~1) is negative and highly significant: $\hat{\beta} = -0.467$ ($\text{SE} = 0.123$, $p < 0.001$), confirming the opportunity-cost result within the four Non-Sahel countries. The Sahel split-sample OLS estimate (Column~2) is also negative in point terms ($\hat{\beta} = -0.877$) but imprecise and insignificant ($\text{SE} = 0.698$, $p = 0.212$), which reflects the smaller sample of 1,866 observations across 130 Sahel districts.

The full-sample interaction model in Column~(3) provides the most informative test. The interaction coefficient is $\hat{\beta}_\text{interaction} = 3.513$ ($\text{SE} = 0.417$, $p < 0.001$), and the baseline (Non-Sahel) yield coefficient is $\hat{\beta} = -1.130$ ($\text{SE} = 0.170$, $p < 0.001$). The net yield effect for Sahel districts is therefore $-1.130 + 3.513 = +2.383$: within Sahel districts, higher predicted yields are associated with \textit{more} conflict in the subsequent 3-month window.

This positive Sahel effect is consistent with the resource-predation channel \citep{DubeVargas2013}: in agro-pastoral environments where armed actors can extract agricultural surplus through coercive grain-market taxation, cattle raiding, and territorial control, a better harvest raises the appropriable surplus and thereby increases the expected payoff to predatory violence. The Sahel is precisely the environment where this mechanism is empirically documented: jihadist and inter-communal armed groups in the Liptako-Gourma region (straddling Burkina Faso, Mali, and Niger) are known to systematically extract resources from agricultural communities \citep{Berman_etal2017}. A larger harvest in this context does not reduce armed-group activity---it increases it by enlarging the prize.

The Sahel Poisson estimate in Column~(5) is negative and significant ($\hat{\beta} = -3.950$, $p < 0.001$), which appears to conflict with the positive interaction result. This divergence is attributable to sample selection in the Poisson absorbing condition, which retains only the 1,806 observations from the 122 Sahel districts that record positive conflict in at least one year. The intensive-margin relationship among persistently conflict-affected Sahel districts---where conflict is a structural feature of the political economy rather than a marginal response to yield variation---may differ from the extensive-margin and full-distribution relationship identified by OLS. The interaction model in Column~(3), which exploits the full distribution and imposes a common fixed-effects structure across both zones, is the preferred basis for the AEZ comparison.

The cross-zone interaction ($p < 0.001$) and the magnitude of the differential effect ($+3.51$ in log-log terms) constitute the most novel empirical contribution of this paper. They demonstrate that aggregate yield--conflict estimates mask a fundamental heterogeneity in the direction of the relationship: the same improvement in agricultural productivity that suppresses conflict in Non-Sahel environments amplifies it in Sahel ones. This finding has direct implications for the design of agricultural interventions in conflict-affected settings, discussed further in Section~\ref{sec:discussion}.

\subsection{Per-Country Results}

Tables~\ref{tab:country_ols} and~\ref{tab:country_poisson} report per-country TWFE OLS log-log and Poisson estimates respectively. Each column is an independent regression for one country.

\begin{table}[H]
\centering
\caption{Per-Country TWFE OLS Log-Log: Effect of Predicted Yield on Conflict}
\label{tab:country_ols}
\begin{threeparttable}

\begingroup
\centering
\begin{tabular}{lccccccc}
   \tabularnewline \midrule \midrule
   Dependent Variable: & \multicolumn{7}{c}{Log conflict (3-mo.)}\\
                             & BFA     & ETH           & MLI          & MWI      & NER     & NGA           & TZA \\   
   Model:                    & (1)     & (2)           & (3)          & (4)      & (5)     & (6)           & (7)\\  
   \midrule
   \emph{Variables}\\
   Log pred.\ yield (kg/ha)  & -1.884  & 1.719$^{***}$ & -0.1738      & 0.0491   & 1.699   & 0.0134        & 0.0652\\   
                             & (1.330) & (0.6443)      & (1.224)      & (0.1003) & (1.485) & (0.2166)      & (0.1481)\\   
   Log population            & 1.787   & -1.675        & -1.215$^{*}$ & -0.0760  & 1.011   & 1.101$^{***}$ & 0.2393$^{***}$\\   
                             & (1.504) & (1.379)       & (0.6602)     & (0.1923) & (1.491) & (0.4114)      & (0.0851)\\   
   \midrule
   \emph{Fixed-effects}\\
   Admin-2 FE                & Yes     & Yes           & Yes          & Yes      & Yes     & Yes           & Yes\\  
   Year FE                   & Yes     & Yes           & Yes          & Yes      & Yes     & Yes           & Yes\\  
   \midrule
   \emph{Fit statistics}\\
   Observations              & 660     & 1,170         & 750          & 3,705    & 456     & 9,016         & 2,126\\  
   R$^2$                     & 0.59739 & 0.36079       & 0.64402      & 0.39575  & 0.56959 & 0.47220       & 0.22126\\  
   Within R$^2$              & 0.01442 & 0.01268       & 0.01169      & 0.00023  & 0.01063 & 0.00347       & 0.00320\\  
   \midrule \midrule
   \multicolumn{8}{l}{\emph{Clustered (Admin-2 FE) standard-errors in parentheses}}\\
   \multicolumn{8}{l}{\emph{Signif. Codes: ***: 0.01, **: 0.05, *: 0.1}}\\
\end{tabular}
 
\par \raggedright 
\textit{Notes:} Clustered standard errors at the Admin-2 level in parentheses. Each column is a separate TWFE regression for one country. All models include Admin-2 and year two-way fixed effects. Country ISO3 codes: BFA = Burkina Faso, ETH = Ethiopia, MLI = Mali, MWI = Malawi, NER = Niger, NGA = Nigeria, TZA = Tanzania. $^{*}$ $p<0.10$, $^{**}$ $p<0.05$, $^{***}$ $p<0.01$.
\par\endgroup



\end{threeparttable}
\end{table}

\begin{table}[H]
\centering
\caption{Per-Country TWFE Poisson: Effect of Predicted Yield on Conflict}
\label{tab:country_poisson}
\begin{threeparttable}

\begingroup
\centering
\begin{tabular}{lccccccc}
   \tabularnewline \midrule \midrule
   Dependent Variable: & \multicolumn{7}{c}{Conflict events (3-mo.)}\\
                             & BFA     & ETH           & MLI            & MWI     & NER     & NGA         & TZA \\   
   Model:                    & (1)     & (2)           & (3)            & (4)     & (5)     & (6)         & (7)\\  
   \midrule
   \emph{Variables}\\
   Log pred.\ yield (kg/ha)  & -2.637  & 0.4879        & -1.176         & -1.678  & -2.160  & 2.375$^{*}$ & 13.62$^{*}$\\   
                             & (1.672) & (4.401)       & (1.369)        & (3.416) & (2.279) & (1.410)     & (8.097)\\   
   Log population            & -5.198  & -9.634$^{**}$ & -7.636$^{***}$ & -4.430  & 2.928   & -0.7018     & 5.316\\   
                             & (3.330) & (4.361)       & (1.676)        & (4.667) & (4.891) & (1.199)     & (5.564)\\   
   \midrule
   \emph{Fixed-effects}\\
   Admin-2 FE                & Yes     & Yes           & Yes            & Yes     & Yes     & Yes         & Yes\\  
   Year FE                   & Yes     & Yes           & Yes            & Yes     & Yes     & Yes         & Yes\\  
   \midrule
   \emph{Fit statistics}\\
   Observations              & 660     & 966           & 735            & 742     & 385     & 6,624       & 788\\  
   Squared Correlation       & 0.78887 & 0.36854       & 0.69855        & 0.57787 & 0.74266 & 0.50469     & 0.31472\\  
   Pseudo R$^2$              & 0.75298 & 0.51113       & 0.67262        & 0.36695 & 0.72365 & 0.49177     & 0.24613\\  
   BIC                       & 2,564.8 & 5,460.7       & 3,116.9        & 1,132.1 & 1,253.6 & 17,518.5    & 1,125.6\\  
   \midrule \midrule
   \multicolumn{8}{l}{\emph{Clustered (Admin-2 FE) standard-errors in parentheses}}\\
   \multicolumn{8}{l}{\emph{Signif. Codes: ***: 0.01, **: 0.05, *: 0.1}}\\
\end{tabular}
 
\par \raggedright 
\textit{Notes:} Clustered standard errors at the Admin-2 level in parentheses. Each column is a separate TWFE regression for one country. All models include Admin-2 and year two-way fixed effects. Country ISO3 codes: BFA = Burkina Faso, ETH = Ethiopia, MLI = Mali, MWI = Malawi, NER = Niger, NGA = Nigeria, TZA = Tanzania. $^{*}$ $p<0.10$, $^{**}$ $p<0.05$, $^{***}$ $p<0.01$.
\par\endgroup



\end{threeparttable}
\end{table}

Across the seven countries, point estimates are negative for Burkina Faso ($\hat{\beta} = -1.884$, $p = 0.164$), Mali ($\hat{\beta} = -0.174$, $p = 0.888$), Malawi ($\hat{\beta} = 0.049$, $p = 0.625$), Nigeria ($\hat{\beta} = 0.013$, $p = 0.951$), and Tanzania ($\hat{\beta} = 0.065$, $p = 0.660$), though none achieves statistical significance at conventional levels. Niger shows a positive but insignificant coefficient ($\hat{\beta} = 1.699$, $p = 0.261$). The absence of significant per-country results is expected given the limited within-country variation once district and year fixed effects are absorbed, and is consistent with the finding that the pooled effect is identified largely from cross-country variation in the yield--conflict relationship.

\textit{Ethiopia.} The exception is Ethiopia, which shows a positive and statistically significant OLS coefficient ($\hat{\beta} = 1.719$, $\text{SE} = 0.644$, $p = 0.009$). Ethiopia is classified in the Non-Sahel zone and would therefore be expected to exhibit a negative yield--conflict relationship consistent with the opportunity-cost channel. The anomalous positive finding is most plausibly explained by the Tigray and Amhara conflicts of 2020--2022, which represent the most severe conflict episode in the sample. These conflicts escalated in agriculturally productive highland zones---districts with high predicted maize yields---and were driven by political and ethnic factors largely unrelated to agricultural conditions. The coincidence of high satellite-derived productivity and high conflict intensity in these specific districts during these years generates a positive within-district yield--conflict correlation in the fixed-effects sample that likely reflects reverse-causal contamination: conflict occurred in productive areas because those areas held strategic political importance, not because higher yields drove conflict. The result illustrates the limits of cross-time within-country identification in the presence of large-scale political crises.
